\section{问题分析}

本题的核心不是单一模型选择，而是按题目四问建立互相衔接的分析框架。问题一先建立可解释的销量基准模型，形成门店、商品和门店--商品三个粒度的基础预测；问题二在商品关联分析基础上进行类别聚合，解决同类零食整合预测；问题三分析外部变量与销量之间的条件统计关联，为综合模型提供特征解释；问题四在门店--商品粒度上融合历史销量、日历、类别和外部变量，并用严格 7 日递推验证检验其是否超过强基准模型。

\begin{figure}[H]
\centering
\includegraphics[width=0.95\textwidth]{technical_route_phase3.png}
\caption{四问关系与技术路线}
\label{fig:technical_route}
\end{figure}

预测评价不能随机切分，因为随机切分会让未来日期信息进入训练过程。本文使用时间顺序切分、日滚动验证和严格 7 日递推窗口验证。日滚动一步验证评价“每天更新真实销量后预测下一天”的能力；严格 7 日递推验证评价“窗口开始时一次性向未来预测 7 天”的能力，两者不能混用。

为避免不同章节的 WAPE 数值被误读，本文将三种验证口径统一列在表 \ref{tab:validation_modes}。全文以严格 7 日递推作为第四问和最终模型评价的主口径；日滚动一步和 7 日窗口总量验证只用于对应问题的方法筛选或补充复核。

\begin{table}[H]
\centering
\caption{三种验证口径的定义与适用范围}
\label{tab:validation_modes}
\resizebox{\textwidth}{!}{
\begin{tabular}{p{3.0cm}p{5.7cm}p{4.7cm}p{3.0cm}}
\toprule
验证口径 & 定义 & 适用范围 & 主要限制 \\
\midrule
日滚动一步验证 & 对验证期每一天，只使用该日期以前的真实历史销量预测当天；进入下一天时，前一天真实销量可作为历史信息更新。 & 问题一、问题二的基础模型筛选；问题三候选方法的预测参考。 & 更接近每日滚动运营，不等同于一次性预测未来 7 天。 \\
7 日窗口总量验证 & 将连续 7 天预测和真实销量分别汇总后计算误差，关注未来 7 天总量偏差。 & 问题一方法复核和聚合层运营可用性说明。 & 汇总会抵消部分日内误差，不能替代细粒度逐日验证。 \\
严格 7 日递推验证 & 每个窗口只使用窗口起点以前的真实历史训练；窗口内第 2 至第 7 天的滞后和滚动特征由前序预测值递推生成。 & 问题四综合模型、显著性检验和最终模型评价主口径。 & 只有 4 个完整验证窗口，统计功效有限。 \\
\bottomrule
\end{tabular}
}
\end{table}

模型选择遵循三条原则：第一，所有结论来自附件数据或明确的模型计算；第二，优先使用本科数学建模竞赛中可解释、可复现、可写入论文的方法；第三，复杂模型只有在误差稳定下降且解释成本可控时才作为主方案。各问题的方法取舍均以相同时间顺序验证下的误差表现和题目解释需求共同决定。
